\section{Generic maximum entropy distributions}\label{sec:genMaxEnt}

In the previous section we have shown that the members of an exponential family are maximum entropy distributions.
However, while their mean parameters constitute the relative interior of the mean polytope, they never lie on faces of lower dimension.
We now exploit the tensor representation study of such faces (see \secref{sec:meanPolytope}) to characterize the maximum entropy distributions in the general case.

\subsection{Main result}

Our main result generalizes the maximum entropy characterization of \theref{the:maxEntropyInterior} to arbitrary mean parameters.

\begin{theorem}[Generic characterization of Maximum Entropy Solutions]
    \label{the:MAINgenMaxEntChar}
    Let $\sstat$ be a statistic and $\basemeasure$ a base measure.
    %% 1. Any member of a face exponential family is maxent
    For any $\weightparamwith$ and $\facein$ the distribution
    \begin{align*}
        \expdistofat{\sstat,\weightparam,\genfacemeasure}{\shortcatvariables}
    \end{align*}
    is a maximum entropy distribution. \\
    %% 2. Find canonical parametrization by the backward map
    For any $\meanparamwith$ the maximum entropy problem has a feasible distribution, if and only if $\meanparamwith\in\genmeanset$.
    If $\meanparamwith\in\genmeanset$ the unique maximum entropy distribution is $\expdistof{\sstat,\weightparam,\genfacemeasure}$ where $\facesymbol=\faceto{\meanparam}$ (see \lemref{lem:faceToMean}) and
    \begin{align*}
        \weightparam
        = \backwardmapwrtof{\sstat,\genfacemeasure}{\meanparam}
    \end{align*}
    where $\genfacemeasure$ is the refined base measure (see \defref{def:faceMeasure}) and $\backwardmapwrt{\sstat,\genfacemeasure}$ is any backward map to the exponential family $\expfamilyof{\sstat,\genfacemeasure}$ (see \defref{def:forwardBackwardMapping}). \\
    %% 3. Represent maxent as a CAN
    For any hypergraph $\graph$ and any face $\facein$ we have that all maximum entropy distributions to mean parameters on $\interiorof{\facesymbol}$ are in $\cansof{\sstat,\graph,\basemeasure}$ if $\facesymbol$ is representable with respect to $\graph$ (see \defref{def:faceRepresentability}).
\end{theorem}

To prepare for the proof of this theorem we first show in an auxiliary lemma that we can reduce the set of feasible distributions in \problemref{prob:maxEntropy}.

\begin{lemma}
    \label{lem:maxEntReduction}
    For any $\meanparamwith\in\genmeanset$ and a face $\facein$ with $\meanparamwith\in\sbinteriorof{\facesymbol}$ the solutions of $\mathrm{P}_{\sstat,\meanparam,\basemeasure}$ and $\mathrm{P}_{\sstat,\meanparam,\genfacemeasure}$ coincide.
\end{lemma}
\begin{proof}
    By \lemref{lem:faceMeasureRepCondition} all feasible distributions are representable by the base measure refined to the face $\facesymbol$.
    We have that any for $\mathrm{P}_{\sstat,\meanparam,\basemeasure}$ feasible distribution $\probwith$ satisfies
    \begin{align*}
        \contraction{\probwith,\genfacemeasure} = 1 \,
    \end{align*}
    and thus $\probwith\in\distmeasuredby{\genfacemeasure}$.
    Conversely, any $\probwith\in\distmeasuredby{\genfacemeasure}$ satisfies
    \begin{align*}
        \contraction{\probwith,\genfacemeasure} = \contraction{\probwith,\basemeasure} = 1
    \end{align*}
    and thus $\probwith\in\distmeasuredby{\basemeasure}$.
    Problem $\mathrm{P}_{\sstat,\meanparam,\basemeasure}$ is thus equal to
    \begin{align*}
        \argmax_{\probwith\in\distmeasuredby{\genfacemeasure}} \sentropyofwrt{\probwith}{\basemeasure}
        \stspace
        \contractionof{\probwith,\sencsstatwith,\basemeasurewith}{\selvariable}
        = \genmeanat{\selvariable}
    \end{align*}
    We further have that any $\probwith\in\distmeasuredby{\genfacemeasure}$
    \begin{align*}
        \sentropyofwrt{\probwith}{\basemeasure} = \sentropyofwrt{\probwith}{\genfacemeasure}
    \end{align*}
    and arrive together with the above equivalence at the claim.
\end{proof}


\begin{proof}[Proof of \theref{the:MAINgenMaxEntChar}]
    % 1.
    To show that for any $\weightparamwith$ and $\facein$ the distribution $\expdistof{\sstat,\weightparam,\genfacemeasure}$ is a maximum entropy distribution, we apply \lemref{lem:maxEntReduction} and \theref{the:maxEntropyInterior}. \\
    % 2.
    If $\meanparamwith\notin\genmeanset$ then \problemref{prob:maxEntropy} is not feasible by definition.
    In case $\meanparamwith\in\genmeanset$ we find $\facein$ with $\facesymbol=\faceto{\meanparam}$ and by \lemref{lem:maxEntReduction} the solution of maximum entropy problem \problemref{prob:maxEntropy} coincides with $\mathrm{P}_{\sstat,\meanparam,\genfacemeasure}$.
    Since by \lemref{lem:faceAsRefinedPolytope} the face is the polytope
    \begin{align*}
        \facesymbol
        =\meansetof{\sstat,\genfacemeasure}
    \end{align*}
    we find a weight parameter
    \begin{align*}
        \weightparam
        = \backwardmapwrtof{\sstat,\genfacemeasure}{\meanparam} \, .
    \end{align*}
    and by \theref{the:maxEntropyInterior} $\expdistof{\sstat,\weightparam,\genfacemeasure}$ is the corresponding maximum entropy distribution. \\
    % 3.
    To show the last claim, it is enough to show that for any face $\facesymbol$ representable by $\graph$ and for any $\weightparamin$ we have
    \begin{align*}
        \expdistof{\sstat,\weightparam,\genfacemeasure}
        \in \cansof{\sstat,\graph,\basemeasure} \, .
    \end{align*}
    To this end, we construct activation cores as follows.
    By \defref{def:faceRepresentability} and \lemref{lem:faceMeasureRepresentation} there is a tensor network $\kcoreof{\graph}$ such that
    \begin{align*}
        \genfacemeasurewith
        = \contractionof{\kcoreof{\graph}\cup\{\basemeasurewith\}}{\shortcatvariables} \, .
    \end{align*}
    For each $\selindexin$ we find an edge $\edge(\selindex)\in\edges$ with $\selindex\in\edge(\selindex)$.
    Now we define a tensor network $\tnetof{\graph}$ by the cores to $\edgein$ as
    \begin{align*}
        \hypercoreofat{\edge}{\catvariableof{\edge}}
        = \contractionof{
            \{\kcoreofat{\edge}{\catvariableof{\edge}}\} \cup \{\softactlegwith\wcols\edge=\edge(\selindex)\}
        }{\catvariableof{\edge}} \, .
    \end{align*}
    Then, we have
    \begin{align*}
        \contractionof{\bencsstatwith,\contractionof{\tnetof{\graph}}{\headvariables}}{\shortcatvariables}
        = \contractionof{\bencsstatwith,\softacttensorwith,\contractionof{\kcoreof{\graph}}{\headvariables}}{\shortcatvariables}
    \end{align*}
    and therefore
    \begin{align*}
        \expdistof{\sstat,\weightparam,\genfacemeasure}
        = \frac{
            \contractionof{\{\bencsstatwith\}\cup\tnetof{\graph}}{\shortcatvariables}
        }{
            \contraction{\{\bencsstatwith,\basemeasurewith\}\cup\tnetof{\graph}}
        } \, .
    \end{align*}
    We conclude $\expdistof{\sstat,\weightparam,\genfacemeasure}\in \cansof{\sstat,\graph,\basemeasure}$.
\end{proof}

\subsection{Family of maximum entropy distributions}


The family of maximum entropy distributions is the set
\begin{align*}
    \maxentfamilyof{\sstat}{\basemeasure} =
    \{\probwith \wcols \exists\meanparam\in\genmeanset : \probwith \quad\text{is a solution of \problemref{prob:maxEntropy}} \} \, .
\end{align*}
\theref{the:MAINgenMaxEntChar} suggests a parametrization of this family as tuples of corresponding faces $\facesymbol$ and weight vectors $\weightparamwith$.
As sketched in \figref{fig:maxEntropyActcore} the activation tensor of each maximum entropy distribution is decomposed into a possibly non-elementary Boolean tensor corresponding to the base measure refinement on a face $\facesymbol$, and an elementary tensor dependent on the weight vector $\weightparamin$.
We thus define the canonical parameter space
\begin{align*}
    \gencanparamset
    \coloneqq \facelatticeof{\genmeanset} \times \weightspace
\end{align*}
and have
\begin{align*}
    \maxentfamilyof{\sstat}{\basemeasure} =
    \left\{
        \frac{
            \contractionof{\softacttensorat{\headvariables},\kcoreofat{\facesymbol}{\headvariables},\bencsstatwith}{\shortcatvariables}
        }{
            \contraction{\softacttensorat{\headvariables},\kcoreofat{\facesymbol}{\headvariables},\bencsstatwith,\basemeasurewith}
        }
        \wcols \canparamwithin,\facesymbol\in\gencanparamset
    \right\} \, .
\end{align*}

Note that these families are disjoint, since each member of an exponential family has support by the support of the face measure.

\begin{figure}[t]
    \begin{center}
        \begin{tikzpicture}[scale=0.35,thick]

    \begin{scope}
    [shift={(-15,0)}]

        \drawtensorblock{2}{-2}{3}{$\probtensor$}
        \drawddwire{0}{-5}{$\catvariableof{0}$}
        \drawddwire{1.5}{-5}{$\catvariableof{1}$}
        \drawddwire{4}{-5}{$\catvariableof{\catorder\shortminus1}$}
        \drawtextnode{3.25}{-4.75}{\colorlabelsize $\cdots$};

        \drawtextnode{8}{-2}{$= \,\frac{1}{\partitionfunction} \cdot $}
    \end{scope}

    \draw[\concolor] (-1,1) rectangle (5,3);
    \drawtextnode{2}{2}{\textcolor{\concolor}{$\kcoreof{\facesymbol}$}}

    \draw[->-] (0,-1)--(0,0);
    \node[right] (text) at (0,0) {\colorlabelsize $\headvariableof{0}$};
    \draw[\concolor] (0,0)--(0,1);
    \drawvariabledot{0}{0}
    \drawtextnode{2}{-0.5}{\colorlabelsize $\cdots$}

    \draw[\probcolor] (0,0) -- (-2,0);
    \draw[\probcolor] (-2,1) rectangle (-4,-1);
    \node[anchor=center,\probcolor] (text) at (-3,0) {\corelabelsize $\softactsymbolof{0,\canparam}$};

    \draw[->-] (4,-1)--(4,0);
    \node[left] (text) at (4,0) {\colorlabelsize $\headvariableof{\seccatorder\shortminus1}$};
    \draw[\concolor] (4,0)--(4,1);
    \drawvariabledot{4}{0}

    \draw[\probcolor] (4,0) -- (6,0);
    \draw[\probcolor] (6,1) rectangle (8.5,-1);
    \node[anchor=center,\probcolor] (text) at (7.25,0) {\corelabelsize $\softactsymbolof{\seccatorder\shortminus1,\canparam}$};

    \drawtensorblock{2}{-2}{3}{$\bencodingof{\sstat}$}
    \drawtdwire{0}{-3}{$\catvariableof{0}$}
    \drawtdwire{1.5}{-3}{$\catvariableof{1}$}
    \drawtdwire{4}{-3}{$\catvariableof{\catorder\shortminus1}$}
    \drawtextnode{3.25}{-4.75}{\colorlabelsize $\cdots$};

\end{tikzpicture}
    \end{center}
    \caption{
        Tensor network decomposition of maximum entropy distributions to the constraint $\meanparamat{\selvariable}=\contractionof{\probwith,\sencsstatwith}{\selvariable}$.
        Blue: Boolean tensor $\kcoreofat{\facesymbol}{\headvariables}$ dependent on the face $\faceto{\meanparam}$, which represents the refinement of the base measure to the face, and is possibly non-elementary.
        Red: Elementary activation tensor $\softacttensorwith=\bigotimes_{\selenumeratorin}\softactlegwith$ to the weight $\weightparamwith= \backwardmapwrtof{\sstat,\genfacemeasure}{\meanparam}$.
    }\label{fig:maxEntropyActcore}
\end{figure}